%! Choosing a Sample: Four Sampling Techniques — Unit 1, Lesson 1.3 — Slides (Beamer)
% PILOT SHAPE (2026-09-06): modeled on AP Statistics lesson 1.4 minus its AP Practice page.
% GRADUAL RELEASE deck: title → targets → warm-up → hook (unresolved) → I DO divider → two frames
% per notes section → WE DO (a live surface, un-answered) → spoken debrief → close & START THE
% HOMEWORK, which is the individual practice. No group activity, no practice launch, no exit
% ticket. There is no spoiler rule: the notes frames ARE the direct instruction.
\documentclass[aspectratio=169, 11pt]{beamer}
\usepackage{saar-beamer}   % provides \ceruleanheader{} and \sectionlabel[color]{}

\newcommand{\redink}[1]{{\color{redacc}\bfseries #1}}

\begin{document}

% TITLE SLIDE (CourseName is not defined in beamer, so the name is literal)
{
\setbeamercolor{background canvas}{bg=cerulean}
\begin{frame}[plain]
  \vspace{2.8cm}\hspace{0.55cm}%
  \begin{minipage}{0.9\textwidth}
    {\color{goldacc}\bfseries\small LESSON 1.3}\\[0.5cm]
    {\color{white}\bfseries\LARGE Choosing a Sample: Four Sampling Techniques}\\[1.2cm]
    {\color{frostmid}\small Statistical Analysis \& Algebraic Reasoning~$\cdot$~Unit 1}
  \end{minipage}
\end{frame}
}

% ── LEARNING TARGETS ─────────────────────────────────────────────────────────
\begin{frame}[t, plain]
  \ceruleanheader{Today's targets}
  \sectionlabel{Lesson 1.3}
  By the end of today, I can\ldots
  \begin{itemize}\setlength{\itemsep}{5pt}
    \item name the \textbf{sampling frame} and say whether \textbf{chance} or
          \textbf{convenience} chose a sample --- by asking \emph{who could never have been picked}.
    \item run a \textbf{simple random sample} from generator digits and a \textbf{systematic
          sample} with $k = N \div n$.
    \item compute a \textbf{stratified} allocation and count what a \textbf{cluster sample} reaches.
    \item tell a \textbf{stratum} from a \textbf{cluster}, and choose --- with a reason --- the
          technique a context calls for.
  \end{itemize}
  \begin{block}{How today works}
    Warm-up $\to$ \textbf{notes together} $\to$ \textbf{one survey we plan as a class} $\to$
    debrief $\to$ \textbf{you start the homework here, alone}.
  \end{block}
\end{frame}

% ── WARM-UP ──────────────────────────────────────────────────────────────────
\begin{frame}[t, plain]
  \ceruleanheader{Warm-Up --- 5 minutes, silent}
  \sectionlabel{Riverbend High School: 900 students, 60 to survey}
  \vspace{0.3cm}
  \begin{center}
    \begin{minipage}{0.86\textwidth}
      \begin{enumerate}\setlength{\itemsep}{7pt}
        \item $270$ of the $900$ students are in 9th grade. What percent of the school is that?
        \item If 9th graders should be the same share of the \emph{survey}, how many of the
              $60$ should be 9th graders?
        \item One list of all $900$, in order. Pick $60$ spread evenly: how many names apart?
              Start at the 7th name --- which two names come next?
      \end{enumerate}
    \end{minipage}
  \end{center}
  \vspace{0.4cm}
  \begin{center}
    {\color{cerulean}\bfseries Algebra 1 arithmetic. Keep the page --- every number on it
    comes back with a name today.}
  \end{center}
\end{frame}

% ── HOOK — leave it UNRESOLVED; the second half of section 2 settles it ──────
{
\setbeamercolor{background canvas}{bg=cerulean}
\begin{frame}[plain]
  \vspace{1.2cm}
  \begin{center}
    {\color{frostmid}\bfseries\small THE COUNCIL NEEDS 60 OF RIVERBEND'S 900 STUDENTS.}\\[0.7cm]
    {\color{white}\bfseries\Large Maya puts the 30 homeroom numbers in a hat, draws two ---
    rooms 4 and 7 --- and surveys all 30 students in each. Pure chance.}\\[0.9cm]
    {\color{goldacc}\bfseries\large Is that as trustworthy as drawing 60 names out of the hat?}\\[0.5cm]
    {\color{frostmid}\small Vote: yes, no, or it depends. Write one line of reason in the hook
    box. We settle this at the end of section 2.}
  \end{center}
\end{frame}
}

% ── I DO — direct instruction begins ─────────────────────────────────────────
{
\setbeamercolor{background canvas}{bg=cerulean}
\begin{frame}[plain]
  \vspace{1.8cm}\hspace{0.8cm}%
  \begin{minipage}{0.86\textwidth}
    {\color{goldacc}\bfseries\small GUIDED NOTES \ \textbullet\ \ I DO}\\[0.7cm]
    {\color{white}\bfseries\Large Open to your Guided Notes page. Fill each blank as we
    name it --- not before.}\\[0.9cm]
    {\color{frostmid}\small Five terms go in the vocabulary box today. Write each one the
    moment we use it.}
  \end{minipage}
\end{frame}
}

% ── 1a. WHO CHOSE THE SAMPLE ─────────────────────────────────────────────────
\begin{frame}[t, plain]
  \ceruleanheader{Who chose the sample?}
  \sectionlabel{notes section 1 $\cdot$ the frame, chance, convenience}
  \vspace{4pt}
  \begin{columns}[t]
    \begin{column}{0.47\textwidth}
      \begin{block}{Probability sample}
        \textbf{Chance} chooses. Every individual on the \textbf{sampling frame} --- the list
        you can actually draw from --- has a known chance of being picked.
      \end{block}
    \end{column}
    \begin{column}{0.47\textwidth}
      \begin{block}{Convenience sample}
        The \textbf{researcher} takes whoever is easiest to reach. Tuesday evening chose the
        pool sample. The bell schedule chose the third-period sample.
      \end{block}
    \end{column}
  \end{columns}
  \vspace{12pt}
  \begin{center}
    {\large \emph{Random} does not mean \emph{unplanned}. Ask: \redink{who could never have
    been picked?}}\\[6pt]
    {\small\color{slate} If that list is not empty, chance did not choose --- and no sample
    size fixes it.}
  \end{center}
\end{frame}

% ── 1b. SRS AND SYSTEMATIC ───────────────────────────────────────────────────
\begin{frame}[t, plain]
  \ceruleanheader{Simple random --- and the systematic shortcut}
  \sectionlabel{notes section 1 $\cdot$ number the frame, then let chance draw}
  \vspace{2pt}
  \begin{block}{Simple random sample: Riverbend's 900 students, numbered 001--900}
    {\ttfamily\large 7 1 4 \quad 0 2 8 \quad 9 5 1 \quad 3 6 6}\\[4pt]
    student 714 \quad$\cdot$\quad student 28 \quad$\cdot$\quad
    \redink{951 --- out of range, skip it} \quad$\cdot$\quad student 366\\[2pt]
    {\small Three digits at a time because $900$ has three. Skip anything out of range or
    already used --- \emph{skip}, not stop.}
  \end{block}
  \vspace{4pt}
  \begin{block}{Systematic sample: the list is already in order}
    $k = N \div n = 900 \div 60 = \mathbf{15}$. Random start between 1 and 15 --- it lands on
    \textbf{7}. Take 7, 22, 37, 52, \dots\ (the warm-up, with its name).\\[2pt]
    {\small Only the \textbf{start} is random. Fine, as long as the list's order has nothing to
    do with the interval.}
  \end{block}
\end{frame}

% ── 2a. STRATIFIED ───────────────────────────────────────────────────────────
\begin{frame}[t, plain]
  \ceruleanheader{Stratified: some from every group}
  \sectionlabel{notes section 2 $\cdot$ strata are alike inside, different between}
  \vspace{2pt}
  \begin{columns}[t]
    \begin{column}{0.46\textwidth}
      \begin{center}
        \small
        \renewcommand{\arraystretch}{1.3}
        \begin{tabular}{@{} l c c @{}}
        \hline
        \textbf{Grade} & \textbf{Students} & \textbf{$1/15$ of each} \\
        \hline
        9th  & 270 & 18 \\
        10th & 240 & 16 \\
        11th & 210 & 14 \\
        12th & 180 & 12 \\
        \hline
        \textbf{Total} & \textbf{900} & \textbf{60} \\
        \hline
        \end{tabular}
      \end{center}
    \end{column}
    \begin{column}{0.5\textwidth}
      \begin{itemize}\setlength{\itemsep}{6pt}
        \item Seniors drive. Grade is a natural set of \textbf{strata}.
        \item $60/900 = 1/15$, so take $1/15$ of \emph{every} grade --- the warm-up's 18 was
              this row.
        \item Stratifying \redink{guarantees} each grade is heard in proportion. A fair SRS only
              makes it \emph{likely}.
      \end{itemize}
    \end{column}
  \end{columns}
\end{frame}

% ── 2b. THE CRUX ─────────────────────────────────────────────────────────────
\begin{frame}[t, plain]
  \ceruleanheader{Cluster: everyone from a few groups --- and Maya's draw}
  \sectionlabel[redacc]{notes section 2 $\cdot$ the whole point of today}
  \vspace{2pt}
  \begin{itemize}\setlength{\itemsep}{4pt}
    \item Cluster sample: draw a \textbf{few whole groups} at random, survey \textbf{everyone}
          in them. $30$ homerooms of $30$; Maya drew rooms $4$ and $7$: $2 \times 30 = 60$.
    \item The fact Maya did not check: \textbf{homerooms 1--9 are 9th grade}, 10--17 are 10th,
          18--24 are 11th, 25--30 are 12th.
    \item Rooms 4 and 7: \redink{60 ninth graders. Zero seniors.} Did chance choose?
          \textbf{Yes.} Did anyone count wrong? \textbf{No.} \textbf{Vote again.}
  \end{itemize}
  \begin{block}{``Chance chose'' is not enough --- the technique has to match the groups}
    A cluster sample works only when each cluster is a \textbf{small mix of the whole}. A
    Riverbend homeroom is \textbf{alike inside} --- it is really a \textbf{stratum} --- and
    drawing a few whole strata can miss a whole grade (about one draw in four lands both rooms
    in one grade). Alike inside $\Rightarrow$ \redink{some from every group}.
  \end{block}
\end{frame}

% ── 2c. THE CONTRAST ─────────────────────────────────────────────────────────
\begin{frame}[t, plain]
  \ceruleanheader{The two everybody mixes up}
  \sectionlabel{notes section 2 $\cdot$ stratified vs.\ cluster}
  \vspace{0.2cm}
  \begin{center}
    \small
    \renewcommand{\arraystretch}{1.35}
    \begin{tabular}{@{} l l l @{}}
    \hline
     & \textbf{Stratified} & \textbf{Cluster} \\
    \hline
    The groups are     & alike inside, different from each other & each a small mix of the whole \\
    You use            & \textbf{all} of the groups & \textbf{a few} of the groups \\
    From each you take & a random sample & \textbf{everyone} \\
    \hline
    \end{tabular}
  \end{center}
  \vspace{0.6cm}
  \begin{center}
    {\color{cerulean}\bfseries Stratified is some from every group.\\[4pt]
     Cluster is everyone from a few groups.}\\[8pt]
    {\small\color{slate} Do not decide from the group names. Decide from \emph{how the groups
    are built}.}
  \end{center}
\end{frame}

% ── WE DO — a LIVE surface: task and data only, no answers ───────────────────
\begin{frame}[t, plain]
  \ceruleanheader{Guided Practice: Cedar Ridge Apartments}
  \sectionlabel[goldacc]{We do $\cdot$ you hold the pen}
  $600$ households in $30$ buildings of $20$; a parking survey of $60$. Buildings 1--6 are
  studios ($120$), 7--20 one-bedrooms ($280$), 21--30 two-bedrooms ($200$). Two-bedroom
  households own the most cars.
  \begin{itemize}\setlength{\itemsep}{4pt}
    \item[(a)] Five plans --- name the technique in each.
    \item[(b)] The stratified plan: $1/10$ of each unit type. Show the one-bedroom row.
    \item[(c)] The hat gives buildings 2, 5, and 6. Who is in the sample --- and is a building a
          mix of the whole, or a stratum?
    \item[(d)] Which plan should management run, \emph{because the groups are}\,\ldots?
    \item[(e)] They have a list of buildings, never of residents. Which techniques survive?
  \end{itemize}
  \begin{block}{The question behind the question}
    If you defended the cluster draw in (c) with ``but it was random,'' you are still at the hook vote.
  \end{block}
\end{frame}

% ── DEBRIEF ──────────────────────────────────────────────────────────────────
\begin{frame}[t, plain]
  \ceruleanheader{Debrief: chance chose --- and it still missed}
  \sectionlabel{10 minutes $\cdot$ pens down, papers up --- we say these aloud}
  \vspace{4pt}
  \begin{enumerate}\setlength{\itemsep}{5pt}
    \item Maya's rooms 4 and 7 back on the board. Say the sentence: \emph{chance chose, nobody
          counted wrong, and not one senior was asked --- because a homeroom is a stratum.}
    \item The third-period class. Who could never have been picked? Then what kind of sample is it?
    \item Cedar Ridge (d): the reason, not the word --- \emph{because the groups are}\,\ldots
  \end{enumerate}
  \vspace{6pt}
  \begin{block}{Cold check --- hands down, thirty seconds}
    Bayside Middle School: $750$ students in three grades, homerooms grouped by grade, $30$ to
    survey. \textbf{Plan 1:} draw $2$ of the $30$ homerooms and ask everyone in them.
    \textbf{Plan 2:} a random $10$ from each grade. \textbf{Name each technique. Which one
    \emph{guarantees} all three grades are heard --- and why doesn't the other, even though it
    is random?}
  \end{block}
\end{frame}

% ── YOU DO — the homework is the individual practice, started in class ───────
\begin{frame}[t, plain]
  \ceruleanheader{You do: start the homework}
  \sectionlabel[goldacc]{Last 10 minutes $\cdot$ on your own $\cdot$ finish it at home}
  \begin{itemize}\setlength{\itemsep}{5pt}
    \item \textbf{Harbor Point Community Pool, done properly} --- the board threw out the
          Tuesday-evening study; you plan $100$ of $1{,}500$.
    \item Run all four plans: generator digits, a stratified allocation, an interval down the
          list, four swim sessions.
    \item Two plans both ``use groups'' --- name which is which, then explain why a random draw
          of lap-swim hours can still miss.
    \item The board's requirement: \emph{seniors must be heard}. Choose the plan and say why.
  \end{itemize}
  \begin{block}{Silent and alone --- I am circulating. Packet pages, scored out of 12, due at the start of the first class after two study halls.}
    Start with \textbf{1, 5, and 6}. Item 5 is the one I am reading over your shoulder ---
    if your reason is ``the draw was unlucky,'' flag me before you leave.
  \end{block}
\end{frame}

% ── CLOSE ────────────────────────────────────────────────────────────────────
{
\setbeamercolor{background canvas}{bg=cerulean}
\begin{frame}[plain]
  \vspace{1.5cm}\hspace{0.8cm}%
  \begin{minipage}{0.86\textwidth}
    {\color{goldacc}\bfseries\small BEFORE YOU GO}\\[0.7cm]
    {\color{white}\bfseries\Large In a probability sample, chance chooses. But chance
    choosing is not enough --- the technique has to match how the groups are built. Alike
    inside: some from every group. A mix of the whole: everyone from a few.}\\[0.9cm]
    {\color{frostmid}\small Next: Lesson 1.4 --- you can choose a sample perfectly and still
    get a misleading answer. Who is missing from the frame, who refuses, and how the question
    is worded: \emph{bias}.}
  \end{minipage}
\end{frame}
}

\end{document}
